% A.115: every relative depth, with an explicit orbit-coinvariant proof.
\subsection{Relative quadratic reciprocity at every depth}

\begin{statement}{A.115}{Proved}
{All relative alternating forms and their simultaneous local values}
Fix $a\in\Q^\times$ with $v_3(a)=1$, put
$K_n=\Q(\alpha_n)$, $\alpha_n^{3^n}=a$, and choose
$\alpha_n^3=\alpha_{n-1}$, with $K_0=\Q$.
For every integer $n\geq1$ set
\[
 \begin{gathered}
 N=3^n,\qquad d=3^{n-1},\qquad
 G=C_N\rtimes(\Z/N)^\times,\\
 T_n=\ker\bigl(\operatorname{Res}_{K_n/\Q}\mathbf G_m
       \xrightarrow{N_{K_n/K_{n-1}}}
                  \operatorname{Res}_{K_{n-1}/\Q}\mathbf G_m\bigr),
 \qquad M_n=T_n[3].
 \end{gathered}
\tag{A.115.1}
\]
The actual action on the $N$ embeddings identifies
$M_n=Z^d\mathbf F_3[Z]/(Z^N)\otimes\chi$, where
$x=1+Z$ is the translation and $\chi$ is the mod-three
cyclotomic character.  In geometric coordinates it is
the kernel of the block-sum map
\[
 (v_i)_{i\in\Z/N}\longmapsto
       (v_h+v_{h+d}+v_{h+2d})_{h\in\Z/d},
\tag{A.115.2}
\]
with its cyclotomic twist.  Put
\[
 B_*(v,w)=\sum_{h=0}^{d-1}
       (v_h w_{h+d}-v_{h+d}w_h),\qquad
                 z=Z^2/(1+Z).
\tag{A.115.3}
\]
The division denotes multiplication by the unit
$(1+Z)^{-1}$ in the stated truncated ring.
The entire space of $G$-equivariant alternating
forms $M_n\times M_n\to\mu_3$ has the following basis:
\[
 \begin{gathered}
 B_j(v,w)=B_*(v,z^{\,d-3^j}w),\qquad 0\leq j<n,\\
 \operatorname{rank}B_j=2\cdot3^j,\qquad
 \operatorname{rad}B_j=Z^{d+2\cdot3^j}
                      \mathbf F_3[Z]/(Z^N)\otimes\chi.
 \end{gathered}
\tag{A.115.4}
\]
Consequently this form space has dimension $n$.
Among its projective forms, rank $2\cdot3^j$
occurs exactly $3^j$ times.

Use $\operatorname{inv}_3$ and the identification
$(\tfrac13\Z)/\Z=\mathbf F_3$ that sends $1/3$ to one.
For the continuous local cohomology groups define
\[
 Q_j(\gamma)=\operatorname{inv}_3
       \bigl(\tfrac12(B_j)_*(\gamma\cup\gamma)\bigr).
\tag{A.115.5}
\]
The actual local Kummer map and these quadratic maps
give a surjection
\[
 T_n(\Q_3)/3T_n(\Q_3)
   \xrightarrow{\ \sim\ }H^1(\Q_3,M_n)
   \xrightarrow{\ (Q_0,\ldots,Q_{n-1})\ }
                         \mathbf F_3^{\,n}.
\tag{A.115.6}
\]
Here multiplication by three on the torus is written
additively in the quotient notation.  Thus all the
displayed quadratic invariants can be prescribed
simultaneously by an actual local relative norm-one
element at every finite depth.
\end{statement}
\begin{proof}
\emph{The module, the form, and the centralizer.}
The embedding and block-norm comparison of A.113
is valid for the fields in (A.115.1): restriction
of the coordinatewise product norm to three-torsion
is exactly (A.115.2).  In the group algebra
\[
 A=\mathbf F_3[C_N]=\mathbf F_3[Z]/(Z^N)
\]
the homomorphism $A\to\mathbf F_3[C_d]$ sends $x$ to
the translation of order $d$.  Its kernel is the
ideal $(x^d-1)=(Z^d)$, proving the displayed module
identification with its actual arithmetic action.
The coefficient field $\Q_3(\mu_N)$ contains no cube
root of $a$: such a root would generate a degree-three
non-Galois field inside an abelian extension, whereas
its Eisenstein cubic has nontrivial quadratic
discriminant class $-3$.
Kummer theory therefore gives degree $N$ over
$\Q_3(\mu_N)$ and the full group $G$ locally.
The same automorphisms give the global group.

On each block the condition $v_h+v_{h+d}+v_{h+2d}=0$
leaves a two-dimensional space.  A cyclic permutation
of its three coordinates has determinant one, and
the permutation of the three coordinates induced by
a unit $u$ has determinant $u\bmod3$.
Translations permute blocks and cyclically rotate
their coordinates.  Thus (A.115.3) is a nondegenerate
$G$-equivariant alternating form with values in
$\mu_3$; the two twists in its arguments multiply
to one and the remaining determinant is $\chi$.

The translation on $M_n$ is one Jordan block of
length $2d$.  An endomorphism commuting with it is
uniquely multiplication by an element of
\[
                         B=A/(Z^{2d}).
\tag{A.115.7}
\]
Conjugation by $u\in H=(\Z/N)^\times$ sends
$x$ to $x^u$.  The twist cancels on endomorphisms,
so $\operatorname{End}_G(M_n)=B^H$.
The adjoint of $x$ for $B_*$ is $x^{-1}$.
Since $-1\in H$, each element of $B^H$ is
self-adjoint.  In characteristic three a
self-adjoint operator $P$ gives the alternating
form $B_*(v,Pw)$: its diagonal is the negative
of itself and hence zero.  Conversely the perfect
pairing $B_*$ writes every equivariant alternating
form in this way.  We have thus constructed the
actual vector-space isomorphism
\[
 B^H\longrightarrow
       \operatorname{Hom}_G(\textstyle\bigwedge^2M_n,\mu_3),
                 \qquad P\longmapsto B_*(\ ,P\ ).
\tag{A.115.8}
\]
\emph{The invariant dimension at every depth.}
The coefficient pairing on $A$,
\[
       \langle f,g\rangle=[x^0]f(x)g(x^{-1}),
\tag{A.115.9}
\]
is perfect and $H$-invariant.  The orthogonal
complement of $(Z^{2d})$ is $(Z^d)$: multiplication
and the involution $Z\mapsto-Z/(1+Z)$ identify the
orthogonal ideal with the annihilator, which is
$(Z^{N-2d})$.
Hence (A.115.9) gives an $H$-equivariant isomorphism
$B^*\simeq I$, where
$I=\ker(\mathbf F_3[C_N]\to\mathbf F_3[C_d])$.
It follows that
\[
                         \dim B^H=\dim I_H.
\tag{A.115.10}
\]
We compute the coinvariants on the right by the
literal orbits of $H$ on the two index sets.
For $0\leq j\leq n-2$, the indices of valuation
$j$ form an orbit of length $2\cdot3^{n-j-1}$
in $\Z/N$ and an orbit of length
$2\cdot3^{n-j-2}$ in $\Z/d$.
Reduction of indices gives the permutation-module
surjection between these two orbits.
Write $m=3^{n-j-1}$.
The decomposition $H=C_2\times C_d$ identifies its
kernel, with the actions of the two indicated
factors, as
\[
 \mathbf F_3[C_2]\otimes
       Y^{m/3}\mathbf F_3[Y]/(Y^m),
\tag{A.115.11}
\]
where a generator of the effective cyclic
three-group acts by $1+Y$.
Its coinvariants have dimension one: the
$C_2$-coinvariants of the first factor have
dimension one, and quotienting the second
factor by multiplication by $Y$ also leaves
exactly one dimension.
The remaining block is $\{0,d,2d\}\to\{0\}$.
Its kernel has basis $e_d-e_0,e_{2d}-e_0$;
the involution swaps them and the three-group
fixes them, so its coinvariants again have
dimension one.
These kernels form a direct sum, because all
their target orbits are disjoint.  There are
$n-1$ kernels of (A.115.11) and one final
kernel.  This proves $\dim B^H=n$, including
$n=1$, when only the final block occurs.

For completeness the asserted invariant basis
is explicit.  The involution fixes
$z=x+x^{-1}-2=Z^2/(1+Z)$, and the unit two sends it to
\[
                              z\longmapsto z+z^2.
\tag{A.115.12}
\]
The element two generates $(\Z/3^n)^\times$:
the binomial expansion gives
$v_3(4^{3^r}-1)=r+1$ for every $r\geq0$,
and reduction modulo three gives its factor
of order two.
For $k=d-3^j>0$, the first nonconstant term
of $(1+z)^k$ has degree $3^j$.  Thus
$(z+z^2)^k-z^k$ is divisible by $z^d=0$
in $B$.
For $j=n-1$ the exponent is zero and the
constant is fixed as well.
The $n$ resulting operators are independent,
since their $Z$-orders $2(d-3^j)$ are
distinct and all less than $2d$.
The dimension calculation proves that they
are a basis of $B^H$.
An operator with first nonzero $Z$-order $r$
has rank $2d-r$ on the one-block module and
kernel its terminal ideal of dimension $r$.
This proves (A.115.4).  In a nonzero linear
combination the largest index $j$ with a
nonzero coefficient determines the smallest
$Z$-order and hence its rank.  Normalize that
coefficient to one; its $j$ preceding
coefficients are free in $\mathbf F_3$.
The projective count is therefore $3^j$.

\emph{The local cohomology maps and all mixed terms.}
For every even $s$, $0\leq s\leq2d$, put
\[
 R(s)=Z^{d+s}A\otimes\chi,
                         \qquad R(2d)=0.
\tag{A.115.13}
\]
The translation coinvariant of a nonzero
$R(s)$ is its first power.  The involution
acts on it as $(-1)^{d+s}(-1)=1$.
There is consequently no $G$-equivariant
quotient $R(s)\to\mu_3$.
Local Tate duality gives
\[
                             H^2(\Q_3,R(s))=0.
\tag{A.115.14}
\]
For $s<t$ both even, the translation-fixed
subspace of $R(s)/R(t)$ is its last power.
The involution acts there as
$(-1)^{d+t-1}(-1)=-1$, so its $H^0$ is zero.
Whenever the quotient has a nondegenerate
alternating form into $\mu_3$, it is its
Cartier dual and its $H^2$ is also zero.
Local duality and the Euler characteristic
formula \cite[I, Corollary 2.3 and Theorem 2.8]{BIR_MilneADT}
therefore make its $H^1$ a vector space of
the coefficient dimension, with a perfect
symmetric cup-product pairing induced by that
alternating form.  These statements apply
to the full continuous local cohomology.

The radical of $B_0$ is $R(2)$.
Equations (A.115.14) and the cohomology
sequence make $H^1(\Q_3,M_n)$ surject onto
the $H^1$ of its two-dimensional
nondegenerate quotient.
A nondegenerate quadratic form over
$\mathbf F_3$ in dimension at least two
represents all three scalars.  This follows
by diagonalization: for any two nonzero
coefficients $a,b$, the values of
$aX^2+bY^2$ already include $0,1,2$.
We can therefore prescribe $Q_0$.

For $1\leq j<n$ set $h=3^{j-1}$ and change
a previously chosen class by the image of
\[
 H^1(\Q_3,R(2h))\longrightarrow H^1(\Q_3,M_n).
\tag{A.115.15}
\]
Every earlier form vanishes with either
argument in $R(2h)$, so these changes
preserve all earlier values, including
their mixed terms.
The radical of $B_j$ restricted to $R(2h)$
is exactly $R(4h)$.
Indeed on the module of length $2d$ one has
$(Z^sM_n)^{\perp_{B_*}}=Z^{2d-s}M_n$,
and the operator defining $B_j$ has
$Z$-order $2d-6h$.
For $w\in Z^{2h}M_n$, the condition that
this operator times $w$ belong to
$Z^{2d-2h}M_n$ is exactly
$w\in Z^{4h}M_n$.
Thus the nondegenerate quotient in this
step is $R(2h)/R(4h)$, of dimension $2h$.

By (A.115.14), its $H^1$ is a quotient of
$H^1(\Q_3,R(2h))$.
Choose an $\mathbf F_3$-linear section of
that surjection and then apply (A.115.15).
For the existing class $\gamma$ and a
correction $\delta$ in its image, the
new value is exactly
\[
 Q_j(\gamma+\delta)=Q_j(\gamma)+Q_j(\delta)
       +\operatorname{inv}_3(B_j)_*(\gamma\cup\delta).
\tag{A.115.16}
\]
On the chosen correction space the
quadratic part is nondegenerate and
the last term is a linear functional.
The perfect bilinear pairing expresses
that functional as pairing with one
vector.  Completing the square therefore
makes (A.115.16) a translate of a
nondegenerate quadratic form plus a
constant, which still represents every
scalar of $\mathbf F_3$.
This prescribes the next value without
altering any preceding value.
Induction through the finite list proves
surjectivity onto $\mathbf F_3^n$.

Finally both number fields in the
relative norm have a unique place
above three.  The local relative
cubic norm is onto by A.105, including
the first cubic when $n=1$.
Shapiro's lemma and Hilbert theorem 90
give
\[
 H^1(\Q_3,T_n)=
   K_{n-1,3}^{\times}/N_{K_{n,3}/K_{n-1,3}}
                         K_{n,3}^{\times}=0.
\tag{A.115.17}
\]
The torus Kummer exact sequence consequently
identifies its actual local quotient with
$H^1(\Q_3,M_n)$, proving (A.115.6).
\end{proof}

The theorem supplies simultaneous local freedom at
every relative depth for the whole displayed space
of quadratic reciprocity forms.  The construction
of a global relative multiplier with the required
auxiliary-prime norm conditions remains an open
step of A.1 and A.5.
The retained exact program
\path{code/relative_forms_all_depths_certificate.py}
checks the polynomial invariant spaces through
$n=7$, the literal module and form maps through
$n=4$, and the finite quadratic value calculation.
The all-depth assertions follow from the orbit
decomposition and the cohomology maps proved above.
